\begin{table}[ht]
\centering
\small
\renewcommand{\arraystretch}{1.10}

\begin{tabular*}{\columnwidth}{@{\extracolsep{\fill}}lcccc@{}}
\toprule
\textbf{Model}
& \textbf{Accuracy}
& \textbf{Macro-F1}
& \textbf{Weighted-F1}
& \textbf{Balanced Acc.} \\
\midrule

Logistic Regression
& \textbf{0.991}
& \textbf{0.989}
& \textbf{0.991}
& \textbf{0.989} \\

XGBoost
& 0.990
& 0.988
& 0.990
& 0.987 \\

Linear SVM
& 0.988
& 0.986
& 0.988
& 0.985 \\

Random Forest
& 0.963
& 0.955
& 0.963
& 0.953 \\

Reference Correlation$^\dagger$
& 0.921
& 0.909
& 0.921
& 0.915 \\

\bottomrule
\end{tabular*}

\vspace{3pt}

\parbox{\columnwidth}{
\footnotesize
$^\dagger$ SingleR-style reference-correlation classifier
(Spearman nearest-centroid mapping), following the template-matching
principle used by reference-mapping methods such as
SingleR \citep{aran2019reference}.
}

\caption{
\textbf{Cell-level} prediction of the disease source of a held-out
patient's cells (6-class), pooled over all 361{,}792 out-of-fold cells
from the patient-level 5-fold protocol. This is a different, easier
statistical task than the patient-level diagnosis reported in
Table~\ref{tab:disease}: cells from the same patient are correlated
observations rather than independent draws, and this metric implicitly
weights patients with more cells more heavily. We report it here only as
a cell-level reference point, not as a patient-level diagnostic result.
}
\label{tab:disease_cell_level}

\end{table}
